\section{Obsluha SD karty}

Pro ukládání a načítání obrazů her (ROM) byl systém vybaven rozhraním pro SD karty. Komunikace
probíhá prostřednictvím rozhraní SPI mikrokontroléru RP2350. Aby byla zajištěna modularita kódu a
snadná správa souborů, je softwarové řešení rozděleno do tří abstraktních vrstev:

\begin{enumerate}
    \item \textbf{Hardwarová vrstva (Platform layer)} zajišťuje nízkoúrovňovou inicializaci SPI
        řadiče a GPIO pinů mikrokontroléru RP2350. Tato část je jedinou platformě závislou vrstvou,
        která přímo manipuluje s registry procesoru.
    \item \textbf{Ovladač SD karty (Storage layer)} implementuje komunikační protokol SD karet v
        režimu SPI. Zajišťuje inicializační sekvenci média, identifikaci typu karty (SDSC/SDHC) a
        především blokové čtení a zápis dat (standardizované sektory o velikosti 512 bytů).
    \item \textbf{I/O rozhraní pro FatFs} slouží jako propojovací most (wrapper) mezi ovladačem
        karty a knihovnou souborového systému.
\end{enumerate}

Pro efektivní práci se soubory a kompatibilitu s osobními počítače je zapotřebí implementovat
souborový systém. Na to byl zvolen souborový systém FAT32, který je standardem pro vystavené
systémy, a byla zvolena knihovna FatFs, která realizuje tento souborový systém.

Díky této vrstvené architektuře se SD karta v aplikaci jeví jako standardní disková jednotka.
Emulátor tak může využívat běžné funkce pro práci se soubory (otevření souboru, sekvenční čtení),
což značně usnadňuje implementaci zavaděče ROM. Podrobná dokumentace inicializačních příkazů a
stavových diagramů SD karty je pro zájemce uvedena v příloze F.

% \section{Obsluha SD karty}
%
% Tato kapitola dokumentuje návrh a implementaci vstupně-výstupní vrstvy určené pro komunikaci se SD
% kartou prostřednictvím sběrnice SPI na mikrokontroléru. Hlavním cílem této části je vytvoření
% spolehlivého abstrakčního rozhraní mezi hardwarovými periferiemi a souborovým systémem
% \textbf{FatFs}, jenž ke svému běhu vyžaduje standardizovaný blokový přístup k úložnému médiu.
%
% Navržené řešení je strukturováno do několika logických úrovní. Tento přístup odděluje hardwarově
% závislé části kódu od vyšší aplikační logiky, čímž je zajištěna modularita a vysoká míra
% přenositelnosti (portability) případné implementace na jiné platformy.
%
% \subsection{Architektura řešení}
%
% Z hlediska softwarové architektury je systém dekomponován do tří základních vrstev:
%
% \begin{itemize}
%     \item Hardwarová vrstva (platform layer)
%     \item Ovladač SD karty (storage layer)
%     \item Disk I/O rozhraní pro FatFs (disk I/O layer)
% \end{itemize}
%
% Tato struktura definuje jasné hranice odpovědnosti pro jednotlivé části systému a významně usnadňuje
% případné budoucí rozšiřování či údržbu kódu.
%
% \subsection{Inicializace SD karty}
%
% Proces inicializace paměťového média probíhá ve standardním režimu SPI a striktně dodržuje sekvenci
% definovanou specifikací fyzické vrstvy SD karet. Po prvotním nastavení sběrnice SPI na nízkou
% přenosovou rychlost je karta uvedena do definovaného výchozího stavu pomocí série hodinových impulsů
% při záměrně deaktivovaném signálu \textbf{CS} (Chip Select).
%
% Následně je opakovaně odesílán příkaz \textbf{CMD0}, jehož účelem je softwarový reset a přepnutí
% karty do stavu \textit{IDLE}. Vzhledem k technologickým prodlevám nemusí být odpověď karty okamžitě
% validní, proto je implementován robustní mechanismus dotazování s omezeným počtem pokusů (timeout).
%
% Po úspěšném přechodu do stavu \textit{IDLE} je prostřednictvím příkazu \textbf{CMD8} ověřena
% kompatibilita karty se specifikací SD verze 2.0 a novější. Validní odpověď formátu R7 zároveň slouží
% k potvrzení podporovaného napěťového rozsahu.
%
% Samotná inicializace média pokračuje specifickou sekvencí příkazů \textbf{CMD55} a \textbf{ACMD41},
% přičemž pro podporu moderních paměťových médií je v argumentu aktivován příznak \textit{HCS} (High
% Capacity Support). Jakmile je inicializační smyčka úspěšně ukončena, provede se vyčtení registru OCR
% pomocí příkazu \textbf{CMD58}. Na základě obsahu tohoto registru je identifikován typ vložené karty
% (\textbf{SDSC} nebo \textbf{SDHC/SDXC}), což je klíčový parametr pro správný výpočet adresace
% fyzických sektorů.
%
% V závěrečné fázi inicializace je komunikační rozhraní SPI přepnuto na maximální podporovanou
% přenosovou rychlost pro efektivní provoz.
%
% \subsection{Hardwarová vrstva (platform layer)}
%
% Hardwarově závislá vrstva plní funkci abstrakce mikrokontroléru. Zajišťuje přímou obsluhu periferií,
% konktrétně SPI řadiče, sdílených datových linek a vyhrazených GPIO pinů. Součástí této vrstvy je
% rovněž správa elementárního časování komunikace.
%
% Poskytované funkční rozhraní zahrnuje:
%
% \begin{itemize}
%     \item Obousměrný přenos jednotlivých bajtů po sběrnici SPI.
%     \item Dynamické řízení přenosové rychlosti hodin (SCLK).
%     \item Ovládání signálu \textbf{Chip Select (CS)}.
%     \item Nízkoúrovňovou inicializaci SPI periferie.
%     \item Generování deterministických časových prodlev (delay).
% \end{itemize}
%
% Tato vrstva jako jediná přímo interaguje s registry mikrokontroléru RP2350 a je plně platformně
% závislá.
%
% \subsection{Ovladač SD karty (storage layer)}
%
% Vrstva ovladače implementuje vlastní komunikační protokol SD karty nad sběrnicí SPI a ke své
% činnosti využívá služby poskytované hardwarovou vrstvou.
%
% Mezi primární zodpovědnosti modulu patří:
%
% \begin{itemize}
%     \item Exekuce inicializační sekvence (\textbf{CMD0}, \textbf{CMD8}, \textbf{ACMD41},
%         \textbf{CMD58}).
%     \item Detekce a správa typu paměťového média (SDSC, SDHC).
%     \item Realizace blokového čtení dat (\textbf{CMD17}).
%     \item Realizace blokového zápisu dat (\textbf{CMD24}).
%     \item Translace logických adres na fyzické adresy sektorů.
%     \item Validace datových tokenů a zpracování stavových kódů (R1, R2, R7 odpovědi).
% \end{itemize}
%
% V kontextu implementovaného systému je SD karta reprezentována jako unifikované blokové zařízení s
% pevnou velikostí alokační jednotky (sektoru) \SI{512}{\byte}.
%
% Během operace čtení odesílá karta – po úspěšném přijetí a zpracování příkazu – datový token
% \texttt{0xFE}, který mikroprocesoru signalizuje připravenost k přesunu dat. Bezprostředně poté
% následuje přenos samotného bloku o velikosti \SI{512}{\byte} a dvou kontrolních bajtů CRC. Tyto
% bajty jsou v rámci této implementace vyčteny a následně zahozeny, neboť striktní CRC kontrola
% komunikace je v režimu SPI specifikací definována jako volitelná.
%
% Během operace zápisu odpovídá karta po přijetí kompletního datového bloku stavovým tokenem (data
% response token), ze kterého lze dekódovat úspěšnost uložení. Interní proces zápisu dat do flash
% paměti karty je následně signalizován držením datové linky v logické nule (busy state) a ukončen
% jejím uvolněním.
%
% \subsection{I/O rozhraní pro FatFs (disk I/O layer)}
%
% Vrstva diskového I/O rozhraní slouží jako propojovací prvek (wrapper) mezi\linebreak nízkoúrovňovým
% ovladačem SD karty a nezávislým modulem souborového systému \textbf{FatFs}. Samotná knihovna FatFs
% je striktně hardwarově agnostická a vyžaduje implementaci sady definovaných blokových operací.
%
% Tato mezivrstva poskytuje a mapuje následující systémová volání:
%
% \begin{itemize}
%     \item \texttt{disk\_initialize()} – provedení inicializační sekvence média.
%     \item \texttt{disk\_status()} – zjištění aktuálního stavu disku (inicializován, chybí, chráněn proti zápisu).
%     \item \texttt{disk\_read()} – čtení $N$ sektorů dat.
%     \item \texttt{disk\_write()} – zápis $N$ sektorů dat.
% \end{itemize}
%
% Uvedené funkce efektivně překládají unifikované požadavky souborového systému na sérii specifických
% příkazů ovladače SD karty. Z pohledu aplikačního rozhraní knihovny FatFs se SD karta
% jeví jako transparentní úložné zařízení a veškeré detaily týkající se sériové SPI komunikace
% zůstávají skryté na nižších vrstvách abstrakce.
